

\chapter{Desarrollo e Implementación}
\label{cap:desarrolloImplementacion}

\section{Representación del contexto de exploración}

Para comprender las decisiones de diseño del modelo de representación interno, es útil partir de un ejemplo
concreto: la estructura del dataset de The Movie Database (TMDb), que fue la primera y principal fuente de datos
utilizada durante el desarrollo. TMDb organiza su información en torno a cinco tipos
de entidades: películas, series de televisión, personas, productoras y cadenas de televisión.
Cada tipo de entidad posee atributos propios --- como el género o la fecha de estreno en el caso
de las películas, o el género y la fecha de nacimiento --- y mantiene
relaciones directas con entidades de otros tipos. Así, una película puede estar producida por
una o varias productoras, contar con actores y miembros de equipo técnico de la colección de
personas, una serie puede estar asociada a una cadena de televisión concreta, etc. O sea, al final tenemos una 
base de datos suficientemente compleja, con atributos específicos al dominio de cada entidad, y posibles relaciones de 
navegación, con aristas que conectan entidades de igual o distinta naturaleza.

Estas relaciones son bidireccionales en el dataset procesado: si una película referencia a una
persona como \textit{Director}, esa persona referencia a su vez a la película bajo la razón
\textit{Director}. Esto permite recorrer el grafo de relaciones en ambas direcciones desde
cualquier punto de partida.

Al analizar esta estructura se infirieron las dos abstracciones fundamentales del modelo
interno. Por un lado, los \textbf{metadatos}: atributos propios de cada entidad, representados
como pares clave-valor, sobre los que el usuario puede establecer condiciones de inclusión o
exclusión. Por otro lado, las \textbf{referencias}: vínculos entre entidades de distintas
colecciones, etiquetados con un tipo de relación (\textit{reason}), que permiten tanto filtrar
por asociación como navegar de una colección a otra siguiendo esos vínculos. Esta distinción
--- metadatos para describir, referencias para relacionar --- es la que articula toda la
representación del contexto de exploración.

    \subsection{La clase \texttt{State}: filtros y contexto acumulado}

La clase \texttt{State} representa el contexto de exploración de un usuario en un momento dado.
Encapsula, por un lado, los filtros activos sobre la colección actual y, por otro, el historial
de transiciones realizadas entre colecciones durante la sesión actual.

Los filtros se organizan en cuatro estructuras \texttt{HashMap} anidadas. Los filtros de metadatos
(\texttt{mfilters}) mapean cada atributo a un conjunto de valores seleccionados, expresando que las
entidades resultantes deben satisfacer al menos uno de dichos valores por atributo. Su contraparte
negativa (\texttt{notMfilters}) excluye las entidades que coincidan con los valores indicados.
Los filtros de referencia (\texttt{rfilters}) operan sobre relaciones entre colecciones: para
cada colección referenciada y tipo de relación, almacenan el conjunto de identificadores de
entidades con los que debe existir vínculo. De igual modo, \texttt{notRfilters} excluye las
entidades que presenten dichos vínculos. Todos estos mapas están indexados por el identificador
de la colección activa, de forma que el mismo objeto \texttt{State} puede acumular filtros
correspondientes a distintas colecciones visitadas a lo largo de la navegación.

Junto a los filtros, \texttt{State} mantiene una lista de objetos \texttt{Link}. Cada
\texttt{Link} registra una transición entre colecciones: almacena la dirección de la transición
(bidireccional), la colección de origen, el tipo de relación seguida y una copia
del estado en el momento de realizar la transición. Esta lista constituye el historial de
navegación y es lo que permite reconstruir, a la hora de ejecutar una consulta, el camino
completo que el usuario ha recorrido entre colecciones.

\texttt{State} implementa \texttt{Serializable} y facilita un constructor de copia que realiza
una copia profunda de todas sus estructuras internas, necesaria para preservar instantáneas
del estado en distintos puntos de la sesión.

    \subsection{\texttt{StateManager}: gestión del ciclo de vida del estado}

\texttt{StateManager} es el punto de entrada a través del cual los controladores interactúan
con el estado de navegación. Se instancia al inicio de la sesión, cuando el usuario selecciona
un dataset y una colección de partida, y se almacena como atributo de la sesión HTTP bajo la
clave \texttt{navState}.

Internamente, \texttt{StateManager} coordina tres estructuras con ciclos de vida distintos.
La primera es el estado activo (\texttt{cur}), que representa el contexto de exploración en
curso y sobre el que se aplican todas las operaciones de filtrado y navegación. La segunda es
una pila de historial (\texttt{historyStack}), una \texttt{Deque} que el
\texttt{NavigationServlet} alimenta invocando \texttt{save()} antes de cada acción, lo que
permite revertir el estado mediante \texttt{restore()}. La tercera es el \texttt{unionSet},
una lista de estados congelados que acumula los contextos que el usuario ha marcado para
combinar mediante unión; al invocar \texttt{union()}, el estado activo se añade al
\texttt{unionSet} y \texttt{cur} se reinicializa con un estado limpio para continuar la
exploración de manera independiente.

Es responsabilidad íntegra de \texttt{StateManager} la gestión de coherencia entre estas tres estructuras. 
Por ejemplo, al ejecutar \texttt{link()} o \texttt{union()}, es
\texttt{StateManager} quien decide cuándo realizar una copia profunda del estado activo
antes de cambiarlo, garantizando que los estados almacenados en el historial y en el
\texttt{unionSet} no se vean afectados por operaciones posteriores. La API pública
de \texttt{StateManager} refleja el vocabulario de acciones disponibles desde el
frontend --- \texttt{addMetadataFilter}, \texttt{removeReferenceFilter}, \texttt{link},
\texttt{goback}, \texttt{save}, \texttt{restore} y \texttt{union} --- delegando en
\texttt{State} las operaciones sobre filtros, y aplicando directamente la lógica de
coordinación cuando la operación afecta a más de una de sus estructuras internas.

\section{Filtrado clásico y relacional}
\label{sec:filtrado}

En TMDb, la exploración de un dataset raramente se agota filtrando por atributos propios de
una colección, por cuenta del tamaño y complejidad de la collección. Un caso de uso habitual es, por ejemplo, 
encontrar películas de un género concreto dirigidas por una persona específica, o excluir series producidas por una 
determinada compañía. El primer tipo de condición opera sobre los metadatos de la entidad; el segundo,
sobre sus relaciones con entidades de otra colección. Esta distinción entre \textit{filtrado
clásico} --- sobre atributos --- y \textit{filtrado relacional} --- sobre vínculos --- motivó el
diseño de dos mecanismos de filtrado independientes y de naturaleza distinta pero que se aplican de forma 
conjunta en cada consulta.

Ambos mecanismos admiten además una variante por la negativa: es posible no solo incluir entidades
que cumplan alguna condición, sino también excluir explícitamente las que la cumplan. La combinación de
inclusiones y exclusiones, tanto sobre metadatos como sobre referencias, define el espacio de
resultados que el usuario está explorando en cada momento.

    \subsection{Construcción dinámica de queries SQL}

El núcleo de la capa de acceso a datos es la traducción del estado de navegación a una query
SQL ejecutable. Esta traducción se asembla de manera incremental: en cada petición, el
\texttt{NavigationDAO} construye la query de forma incremental añadiendo cláusulas
\texttt{AND} sobre una consulta base que selecciona entidades de la colección activa.

El \texttt{appendFilterClauses} es el método central de este proceso, que recibe el  
\texttt{StringBuilder} de la query en construcción, la lista de parámetros y el estado
actual, y le añade todas las condiciones derivadas de los filtros acumulados. Para el caso
más simple --- sin historial de navegación entre colecciones --- delega directamente en
\texttt{appendStateFilters}, que itera sobre los cuatro mapas de filtros del estado
(\texttt{mfilters}, \texttt{notMfilters}, \texttt{rfilters}, \texttt{notRfilters}) y
genera una subconsulta \texttt{EXISTS} o \texttt{NOT EXISTS} por cada condición activa.
Cuando el estado contiene además un historial de \textit{links}, \texttt{appendFilterClauses}
actúa de forma recursiva, generando subconsultas anidadas que incorporan el contexto de
cada transición previa; este mecanismo se describe en detalle en la
sección~\ref{sec:navegacion_transversal}.

Todas las queries se construyen con parámetros posicionales (\texttt{?}), nunca
interpolando valores directamente en el SQL. Los valores se acumulan en una lista
\texttt{params} y se asignan posteriormente mediante \texttt{stmt.setObject()}, lo que
delega en el \texttt{PreparedStatement} de JDBC la responsabilidad de escapar los valores
antes de enviarlos al motor, para prevenir inyecciones SQL. El coste de filtrado se traslada
íntegramente a PostgreSQL, sin materializar subconjuntos intermedios de resultados en
memoria.

    \subsection{Filtros de metadatos}

Los filtros de metadatos permiten restringir los resultados de una colección en función de los
atributos clave-valor asociados a sus entidades. En el contexto de TMDb, por ejemplo, esto equivale a
operaciones como seleccionar únicamente películas del género \textit{Action}, o excluir
aquellas con idioma original distinto del inglés.

Cada filtro de inclusión se traduce en una subconsulta \texttt{EXISTS} sobre la tabla
\texttt{metadata}. Si el usuario ha seleccionado múltiples valores para el mismo atributo,
se genera una cláusula \texttt{EXISTS} independiente por cada valor, lo que impone que la
entidad deba satisfacer \textit{todos} ellos simultáneamente. Los filtros de exclusión
(\texttt{notMfilters}) siguen la misma estructura pero con \texttt{NOT EXISTS}:

\begin{lstlisting}[language=Java]
for (var entry : state.getMfilters().entrySet()) {
    String attr = entry.getKey();
    for (String val : entry.getValue()) {
        sql.append(" AND EXISTS (SELECT 1 FROM metadata " + m
            + " WHERE " + m + ".entity_id = " + e + ".global_id"
            + " AND " + m + ".key = ? AND " + m + ".value = ?) ");
        params.add(attr);
        params.add(val);
    }
}
\end{lstlisting}

Esta elección implica una semántica de conjunción estricta entre valores del mismo atributo,
lo cual puede parecer contraintuitivo cuando los valores son mutuamente excluyentes
--- como ocurre con \texttt{Genres} en TMDb, donde una película puede pertenecer a varios
géneros a la vez. En ese caso, seleccionar \textit{Action} y \textit{Drama} devuelve
las películas que tienen \textit{ambos} géneros, no las que tienen alguno de los dos.
Esta semántica se mantuvo de forma deliberada para garantizar consistencia con el
comportamiento del resto de filtros del sistema, y la describiremos como parte de los
obstáculos encontrados en la sección~\ref{sec:obstaculos}.

    \subsection{Filtros de referencia}

Los filtros de referencia permiten restringir los resultados de una colección en función
de sus vínculos con entidades concretas de otra colección. En TMDb, esto cubre casos como
mostrar únicamente las películas en las que una persona determinada participó como
\textit{Director}, o excluir series emitidas por una cadena específica.

Cada filtro de referencia queda identificado por tres elementos: la colección referenciada,
el tipo de relación (\textit{reason}) y el identificador de la entidad concreta. Al igual
que con los metadatos, se genera una subconsulta \texttt{EXISTS} por cada combinación
activa, realizando un \texttt{JOIN} entre la tabla \texttt{reference}, la tabla
\texttt{entities} y la tabla \texttt{collections} para verificar la existencia del vínculo:

\begin{lstlisting}[language=Java]
sql.append(" AND EXISTS ("
    + "SELECT 1 FROM reference " + r
    + " JOIN entities " + re + " ON " + r + ".target_id = " + re + ".global_id"
    + " JOIN collections " + rc + " ON " + re + ".collection_global_id = "
    + rc + ".global_id"
    + " WHERE " + r + ".source_id = " + e + ".global_id"
    + " AND " + rc + ".original_id = ?"
    + " AND " + rc + ".dataset_id = ?"
    + " AND " + r + ".reason = ?"
    + " AND " + re + ".original_id = ?) ");
\end{lstlisting}

Los filtros de exclusión (\texttt{notRfilters}) utilizan \texttt{NOT EXISTS} con la misma
estructura. Todos los filtros activos --- de metadatos e de referencia, en sus variantes
positiva y negativa --- se aplican conjuntamente sobre la misma consulta base, de forma
que el resultado final es la intersección de todas las condiciones establecidas por el
usuario en ese momento.

    \subsection{Facetas: cómputo dinámico de valores disponibles}

Las facetas son el mecanismo por el que el sistema informa al usuario de qué valores
siguen siendo relevantes dado el contexto de filtrado actual: qué géneros tienen al
menos una película entre los resultados visibles, qué directores aparecen en esas
películas, o a qué cadenas están asociadas las series que quedan tras aplicar los filtros.
Este cómputo es siempre dinámico: se recalcula en cada petición a partir del conjunto
de entidades que satisfacen el estado actual.

El sistema distingue dos tipos de facetas, servidas conjuntamente por \texttt{FacetsServlet}
a través del endpoint \texttt{GET /api/facets}:

Las \textbf{facetas de metadatos} se obtienen mediante \texttt{getFacets()}, que construye
una subconsulta con los identificadores de las entidades actualmente visibles
--- aplicando todos los filtros del estado --- y sobre ella ejecuta un \texttt{GROUP BY}
sobre la tabla \texttt{metadata}:

\begin{lstlisting}[language=Java, basicstyle=\small\ttfamily]
String sql =
    "SELECT m.key, m.value, COUNT(*) as cnt "
    + "FROM metadata m "
    + "WHERE m.entity_id IN (" + subquery + ") "
    + "GROUP BY m.key, m.value "
    + "ORDER BY m.key, cnt DESC";
\end{lstlisting}

El resultado es un mapa de atributo a lista de valores con su frecuencia, ordenados
de mayor a menor dentro de cada atributo.

Las \textbf{facetas de referencia} se obtienen mediante \texttt{getReferenceFacets()},
que sobre la misma subconsulta de entidades visibles realiza un \texttt{JOIN} con la
tabla \texttt{reference} para descubrir qué entidades de otras colecciones están
vinculadas a los resultados actuales, agrupadas por colección y tipo de relación.
El resultado expone, por ejemplo, qué personas concretas aparecen como \textit{Director}
entre las películas visibles, junto con el número de películas en las que aparecen.

Junto a las facetas, el mismo endpoint devuelve los \textit{links} disponibles
--- obtenidos mediante \texttt{getAvailableLinks()} --- que indican a qué otras
colecciones es posible navegar desde el conjunto actual, y los filtros activos en ese
momento (\texttt{activeFilters}), permitiendo al frontend reflejar el estado de
selección sin necesidad de mantener dicha información en el cliente.

\section{Navegación transversal}
\label{sec:navegacion_transversal}

Hasta ahora, los mecanismos de filtrado descritos operan dentro de una única colección:
dadas las entidades de \textit{Movies}, se restringe el conjunto según atributos propios
o según vínculos con entidades concretas de otras colecciones. Sin embargo, el sistema
permite también \textit{trasladarse} de una colección a otra siguiendo una relación,
arrastrando el contexto de filtrado acumulado. Este mecanismo es lo que se denomina
navegación transversal.

Un ejemplo concreto en TMDb: el usuario está explorando películas con el filtro de género
\textit{Action} activo. Desde ese contexto, decide navegar hacia la colección
\textit{People} siguiendo la relación \textit{Actor}, para ver qué actores aparecen en
esas películas de acción. Tras el salto, la colección activa pasa a ser \textit{People},
pero el sistema debe recordar que solo interesan las personas vinculadas al subconjunto
de películas que cumplían el filtro previo. El historial de transiciones es, por tanto,
parte del contexto de exploración, no un simple registro de navegación.


    \subsection{El mecanismo de \texttt{link} y \texttt{goback}}

La transición entre colecciones se realiza mediante la operación \texttt{link}, que recibe
el identificador de la colección destino y el tipo de relación a seguir. Su implementación
en \texttt{StateManager} tiene tres efectos:

\begin{lstlisting}[language=Java]
public void link(int env, String reason) {
    this.linkList.add(new Link(true, this.cur.getCurrentCollectionId(),
                               reason, new State(this.cur)));
    this.cur.link(env, reason);
    this.cur.setLinks(new ArrayList<>(this.linkList));
}
\end{lstlisting}

Primero, se añade un nuevo objeto \texttt{Link} al \texttt{linkList}, marcado como
transición hacia delante (\texttt{forward = true}). Este objeto almacena la colección de
origen, el tipo de relación seguida y una copia profunda del estado activo en el momento
de la transición, preservando así los filtros que estaban activos antes del salto.
Segundo, se actualiza la colección activa en \texttt{cur} a la colección destino.
Tercero, se sincroniza la lista de links dentro del propio \texttt{State}, de forma que
el \texttt{NavigationDAO} pueda accederla directamente al construir las queries.

La operación \texttt{goback} invierte la última transición registrada en el
\texttt{linkList}. En lugar de eliminar el último \texttt{Link}, añade uno nuevo con la
dirección invertida respecto al anterior ---si el último era hacia delante, el nuevo es
hacia atrás, y viceversa--- y restaura como colección activa la colección de origen
almacenada en ese \texttt{Link}:

\begin{lstlisting}[language=Java]
public void goback() {
    Link last = this.linkList.get(this.linkList.size() - 1);
    if (last.isForward()) {
        this.linkList.add(new Link(false, this.cur.getCurrentCollectionId(),
                                   last.getReason(), new State(this.cur)));
    } else {
        this.linkList.add(new Link(true, this.cur.getCurrentCollectionId(),
                                   last.getReason(), new State(this.cur)));
    }
    this.cur.goback(last.getEnv());
    this.cur.setLinks(new ArrayList<>(this.linkList));
}
\end{lstlisting}

El hecho de que \texttt{goback} añada un nuevo \texttt{Link} en lugar de eliminar el
último tiene una consecuencia relevante: el historial es siempre creciente y nunca se
borra. Esto significa que la query SQL resultante debe recorrer toda la cadena de
transiciones acumuladas, incluyendo los retrocesos, para reconstruir fielmente el
contexto de exploración.


    \subsection{Subconsultas anidadas: traducción del historial de \textit{links} a SQL}

Cuando el estado contiene un historial de transiciones, \texttt{appendFilterClauses}
no puede limitarse a aplicar los filtros del estado activo: debe también expresar en SQL
la restricción que impone cada transición previa. Para ello utiliza una sobrecarga
recursiva que recorre el \texttt{linkList} desde el último \texttt{Link} hacia el primero,
generando en cada nivel una subconsulta anidada.

La estructura de la recursión es la siguiente: para cada nivel de profundidad, se aplican
primero los filtros del estado en ese nivel mediante \texttt{appendStateFilters}, y a
continuación se genera una cláusula \texttt{AND ... IN (SELECT ...)} que restringe las
entidades del nivel actual a aquellas que, a través de la relación registrada en el
\texttt{Link}, están vinculadas a entidades del nivel siguiente. La dirección del
\texttt{Link} determina qué columna de la tabla \texttt{reference} actúa como origen y
cuál como destino:

\begin{lstlisting}[language=Java]
if (link.isForward()) {
    sql.append(" AND " + e + ".global_id IN ("
        + "SELECT " + r_link + ".target_id FROM reference " + r_link
        + " JOIN entities " + next_e + " ON " + r_link + ".source_id = "
        + next_e + ".global_id"
        + " JOIN collections " + next_c + " ON " + next_e
        + ".collection_global_id = " + next_c + ".global_id"
        + " WHERE " + next_c + ".dataset_id = ? AND " + next_c
        + ".original_id = ? AND " + r_link + ".reason = ? ");
} else {
    sql.append(" AND " + e + ".global_id IN ("
        + "SELECT " + r_link + ".source_id FROM reference " + r_link
        + " JOIN entities " + next_e + " ON " + r_link + ".target_id = "
        + next_e + ".global_id"
        + " JOIN collections " + next_c + " ON " + next_e
        + ".collection_global_id = " + next_c + ".global_id"
        + " WHERE " + next_c + ".dataset_id = ? AND " + next_c
        + ".original_id = ? AND " + r_link + ".reason = ? ");
}
appendFilterClauses(sql, params, link.getState(), linkList, index - 1, depth + 1);
sql.append(") ");
\end{lstlisting}

Cada llamada recursiva abre una subconsulta que se cierra después de que la llamada
interna haya completado su propio nivel, resultando en una estructura de paréntesis
anidados que refleja fielmente la cadena de transiciones del usuario. Los alias de tabla
(\texttt{e}, \texttt{e1}, \texttt{e2}, \texttt{rl0}, \texttt{rl1}\ldots) se generan
dinámicamente a partir del parámetro \texttt{depth} para evitar colisiones entre niveles.

Retomando el ejemplo anterior ---películas de acción, navegación hacia \textit{People}
por la relación \textit{Actor}--- la query resultante para obtener los actores visibles
tendría la forma:

\begin{lstlisting}[language=SQL]
SELECT DISTINCT e.original_id, e.name
FROM entities e
JOIN collections c ON e.collection_global_id = c.global_id
WHERE c.dataset_id = ? AND c.original_id = ?   -- People
AND e.global_id IN (
    SELECT rl0.source_id FROM reference rl0
    JOIN entities e1 ON rl0.target_id = e1.global_id
    JOIN collections c1 ON e1.collection_global_id = c1.global_id
    WHERE c1.dataset_id = ? AND c1.original_id = ?  -- Movies
    AND rl0.reason = 'Actor'
    AND EXISTS (SELECT 1 FROM metadata m0         -- Genres = Action
                WHERE m0.entity_id = e1.global_id
                AND m0.key = ? AND m0.value = ?)
)
\end{lstlisting}

La subconsulta interior opera sobre la colección \textit{Movies} con sus filtros activos,
y la consulta exterior recupera únicamente las personas de \textit{People} que aparecen
como origen de una relación \textit{Actor} hacia alguna de esas películas.

\section{Operaciones de conjuntos}
\label{sec:operaciones_conjuntos}

Los mecanismos de filtrado y navegación transversal descritos hasta ahora operan siempre
sobre un único contexto de exploración activo. Sin embargo, hay casos de uso que requieren
combinar los resultados de dos contextos construidos de forma independiente. Un ejemplo
concreto en TMDb: el usuario quiere obtener todas las personas que han trabajado como
\textit{Director} en películas de acción \textit{o} como \textit{Actor} en series de
ciencia ficción. Ninguno de los dos conjuntos es subconjunto del otro, y construir ambas
condiciones simultáneamente dentro de un único estado no es posible con la semántica
conjuntiva del filtrado descrito en la sección~\ref{sec:filtrado}.

Para cubrir este caso, el sistema incorpora una operación de unión: el usuario puede
congelar el contexto actual, iniciar una nueva exploración desde cero y, al consultar
los resultados, obtener las entidades que satisfacen cualquiera de los contextos
acumulados.


\subsection{Construcción del \texttt{unionSet}}

La operación \texttt{union} en \texttt{StateManager} tiene cuatro efectos sobre el estado
interno:

\begin{lstlisting}[language=Java]
public void union(int collectionId) {
    this.cur.setLinks(new ArrayList<>(this.linkList));
    this.unionSet.add(this.cur);
    this.historyStack.clear();
    this.linkList.clear();
    this.cur = new State(this.datasetId, collectionId);
}
\end{lstlisting}

Primero, sincroniza el historial de \textit{links} dentro del estado activo, de forma que
el estado congelado que se va a preservar tenga acceso a la cadena de transiciones
completa. Segundo, añade el estado activo actual al \texttt{unionSet} sin realizar ninguna
copia: dado que \texttt{cur} va a ser reemplazado inmediatamente, no hay riesgo de que
futuras modificaciones lo alteren. Tercero, vacía tanto la \texttt{historyStack} como el
\texttt{linkList}, desvinculando completamente el nuevo contexto del anterior. Cuarto,
inicializa \texttt{cur} con un estado limpio sobre la colección indicada, desde el que el
usuario puede comenzar una exploración independiente.

El \texttt{unionSet} puede acumular un número arbitrario de entradas repitiendo esta
operación. Cada entrada es un estado completo e independiente, con su propia colección
activa, sus propios filtros y su propio historial de transiciones.


\subsection{Resolución de la unión en base de datos}

La consulta del \texttt{unionSet} se realiza a través del endpoint
\texttt{GET /api/union}, servido por \texttt{UnionServlet}, que delega en los métodos
\texttt{getUnionEntities()} y \texttt{countUnionEntities()} del \texttt{NavigationDAO}.

Ambos métodos construyen una única query que itera sobre todas las entradas del
\texttt{unionSet} y concatena un bloque de condiciones por entrada, separados por
\texttt{OR}:

\begin{lstlisting}[language=Java]
for (int i = 0; i < unionSet.size(); i++) {
    State state = unionSet.get(i);
    if (i > 0) sql.append(" OR ");
    sql.append(" (c.dataset_id = ? AND c.original_id = ? ");
    params.add(state.getDatasetId());
    params.add(state.getCurrentCollectionId());
    appendFilterClauses(sql, params, state);
    sql.append(" ) ");
}
\end{lstlisting}

Cada bloque es estructuralmente idéntico a la query que se generaría para ese estado en
una consulta normal: aplica \texttt{appendFilterClauses} sobre el estado correspondiente,
lo que incluye sus filtros de metadatos, sus filtros de referencia y, si procede, las
subconsultas anidadas derivadas de su historial de transiciones. La unión entre bloques
se expresa con \texttt{OR} a nivel de la cláusula \texttt{WHERE}, y el \texttt{DISTINCT}
en el \texttt{SELECT} garantiza que las entidades presentes en más de un bloque no
aparezcan duplicadas en el resultado.

Una consecuencia de este diseño es que la unión puede combinar entradas sobre
colecciones distintas. En el ejemplo anterior, el primer bloque opera sobre
\textit{People} filtrado por películas de acción con relación \textit{Director}, y el
segundo también sobre \textit{People} pero filtrado por series de ciencia ficción con
relación \textit{Actor}. Que ambas entradas compartan colección activa no es un requisito
del sistema: el campo \texttt{collection\_id} devuelto por \texttt{getUnionEntities()}
indica a qué colección pertenece cada entidad del resultado, permitiendo al frontend
distinguirlas si fuera necesario.

\section{Obstáculos encontrados y soluciones implementadas}
\label{sec:obstaculos}